% !TeX program = pdflatex
% Resumo do Consensus Statement "Guidelines for T cell nomenclature"
% Masopust et al., Nature Reviews Immunology 26:298-313 (April 2026)
% DOI: 10.1038/s41577-025-01238-2
% Ênfase: subpopulações (subsets) de células T CD4+ e nomenclatura modular v1.0
\documentclass[11pt,a4paper]{article}

\usepackage[T1]{fontenc}
\usepackage[utf8]{inputenc}
\usepackage{lmodern}
\usepackage[brazil]{babel}
\usepackage[margin=2.2cm]{geometry}
\usepackage{booktabs}
\usepackage{longtable}
\usepackage{array}
\usepackage{multirow}
\usepackage{ragged2e}
\usepackage{xcolor}
\usepackage{enumitem}
\usepackage[hidelinks]{hyperref}

% ----- Macros auxiliares -----
\newcommand{\up}[1]{\textsuperscript{#1}}          % sobrescrito de marcador (+, --, hi, low)
\newcommand{\tc}[1]{T\textsubscript{#1}}           % nome de subset com subscrito: \tc{CM}, \tc{H}1, \tc{reg}
\newcommand{\gd}{\ensuremath{\gamma\delta}}        % gama-delta
\newcommand{\ifng}{IFN\ensuremath{\gamma}}
\newcommand{\rort}{ROR\ensuremath{\gamma}t}

% Coluna de parágrafo justificado à esquerda para tabelas
\newcolumntype{L}[1]{>{\RaggedRight\arraybackslash}p{#1}}

\definecolor{cd4blue}{RGB}{20,80,150}
\definecolor{softgray}{RGB}{95,95,95}

\hypersetup{
  pdftitle={Diretrizes para a Nomenclatura de Celulas T -- Resumo},
  pdfauthor={Resumo automatizado}
}

\title{\textbf{Diretrizes para a Nomenclatura de Células T}\\[4pt]
\large Resumo do \emph{Consensus Statement} com ênfase nas subpopulações de células T CD4\up{+}}
\author{}
\date{}

\begin{document}
\maketitle
\vspace{-2.2em}

\begin{center}
\small\color{softgray}
Baseado em: Masopust D., Awasthi A., Bosselut R.\ \emph{et al.}
``Guidelines for T cell nomenclature'' (Consensus Statement).\\
\emph{Nature Reviews Immunology} \textbf{26}, 298--313 (abril 2026).
\href{https://doi.org/10.1038/s41577-025-01238-2}{doi:10.1038/s41577-025-01238-2}.
\end{center}

\vspace{0.5em}
\hrule
\vspace{1em}

\section*{Resumo}
Os avanços na biologia de células T revelaram uma heterogeneidade que a nomenclatura
tradicional (baseada em poucos ``subsets'' idealizados) não consegue capturar. Isso levou a
um alargamento \emph{ad hoc} das definições e à criação de designações novas e pouco
padronizadas. O documento de consenso propõe três soluções: (1) que artigos primários
\textbf{definam explicitamente, na seção de Métodos}, o critério experimental usado para
designar cada subset; (2) um conjunto de \textbf{definições padronizadas} para os termos
já existentes (Tabelas~\ref{tab:dif}--\ref{tab:mod}); e (3) uma \textbf{nomenclatura
modular v1.0}, que descreve propriedades biológicas individuais de uma população por meio
de descritores curtos e combináveis, em vez de forçar a célula para dentro de um subset
único. Este resumo cobre todo o artigo, mas detalha em especial as subpopulações de
células T CD4\up{+}.

\tableofcontents
\vspace{1em}
\hrule

% =====================================================================
\section{Contexto e motivação}
\begin{itemize}[leftmargin=1.4em,itemsep=2pt]
  \item \textbf{Histórico:} anos 1960 --- linfócitos separados em B e T; anos 1970 ---
        T divididas em CD4\up{+} e CD8\up{+} (co-receptor do TCR); anos 1990 --- a citometria
        de 4 cores populariza os termos memória central (\tc{CM}) e memória efetora (\tc{EM}).
  \item \textbf{O problema:} propriedades de migração, função, proliferação, potencial de
        diferenciação e longevidade \emph{nem sempre são co-reguladas}. Um vocabulário
        limitado de subsets falha, por princípio, em capturar todas essas propriedades.
  \item \textbf{Agravante:} tecnologias de célula única de alta dimensão expõem uma
        heterogeneidade massiva, tornando insustentável a classificação em poucos subsets
        ``tidy''.
  \item \textbf{Origem do consenso:} \emph{workshop} no Keystone Symposia sobre memória de
        células T (Banff, Alberta, fev.\ 2020); recomendações desenvolvidas entre 2020 e 2025,
        lideradas por D.~Masopust e R.~Ahmed.
\end{itemize}

\subsection{As três recomendações (Box 1)}
\begin{enumerate}[leftmargin=1.6em,itemsep=2pt]
  \item Autores devem declarar claramente, sob um título próprio nos \textbf{Métodos}, como
        os subsets relevantes foram definidos --- requisito a ser cobrado por editores e
        revisores.
  \item As definições atuais dos principais subsets ficaram ``amorfas''; é preciso um
        \textbf{dicionário comum} de definições e ensaios/marcadores substitutos (\emph{proxies}).
  \item Deve-se desenvolver um sistema de classificação que \textbf{acomode} (em vez de ser
        rompido por) perfis de alta dimensão, adaptável a novas descobertas e ciente das
        desigualdades técnicas entre áreas, espécies e laboratórios.
\end{enumerate}

\subsection{Definir os termos nos Métodos (Box 2)}
A intenção não é impor regras rígidas, mas dar \emph{transparência}. Exemplos citados:
``\tc{CM} refere-se a células T CD62L\up{+}CD44\up{+}'' ou ``células \tc{H}17 foram definidas por
marcação intracelular de IL-17A após estímulo com peptídeo cognato e expressão de \rort{}''.

% =====================================================================
\section{Estados de diferenciação principais (TCR\texorpdfstring{$\alpha\beta$}{ab}\up{+})}
Definições da Tabela~1 do artigo, aplicáveis tanto a CD4\up{+} quanto a CD8\up{+}.

\renewcommand{\arraystretch}{1.25}
\small
\begin{longtable}{L{2.4cm} L{6.2cm} L{6.6cm}}
\caption{Estados de diferenciação de células T (proxies representativos).}\label{tab:dif}\\
\toprule
\textbf{Estado} & \textbf{Marcadores / ensaios} & \textbf{Descrição} \\
\midrule
\endfirsthead
\toprule
\textbf{Estado} & \textbf{Marcadores / ensaios} & \textbf{Descrição} \\
\midrule
\endhead
\bottomrule
\endfoot
\textbf{Naive (N)} &
Humano: CCR7\up{hi}, CD45RA\up{+}, CD45RO\up{--}, CD95\up{--}.
Camund.: CD62L\up{+}, CD44\up{low}, CD11a\up{low}, CD122\up{low} &
T madura que saiu do timo, nunca estimulada por antígeno cognato na periferia; permanece
indiferenciada. Expressa mal \ifng{}/IL-4/IL-5 em estímulo curto (5--12\,h). \\
\textbf{Efetora} &
Qualquer molécula efetora (ex.: granzima~B, perforina, citocinas antimicrobianas, CD40L) &
Maior expressão de moléculas/genes ``efetores'', típica de células recém-estimuladas.
O termo é ambíguo (implica ou não diferenciação terminal / vida curta). \\
\textbf{Ativada (A)} &
Transitoriamente CD69\up{+}, CD25\up{hi}, KLF2\up{low}; muitas vezes Ki67\up{+}, HLA-DR\up{+},
CD38\up{hi} (humano) &
Estimulação recente via TCR (salvo ``bystander''). Não implica destino, longevidade nem
função específica. \\
\textbf{Memória (M)} &
Humano: CD11a\up{hi}, freq.\ CD95\up{+}, CD58\up{+}, CD49d\up{+}.
Camund.: CD44\up{hi}, CD11a\up{hi}, tipic.\ CD127\up{+} &
Progênie proliferada de uma célula naive que sofreu ativação e diferenciação. Longeva
(meses a décadas), relativamente quiescente e \emph{fenotipicamente diversa}. \\
\textbf{Anérgica (G)} &
Camund.: CD44\up{+}, FR4\up{+}, CD73\up{+}, FOXP3\up{--}, CD304/NRP1\up{+}, CXCR5\up{--}
(humano: mal caracterizado) &
Célula previamente ativada, hiporresponsiva; perde capacidade de produzir IL-2 ao sinal do
TCR (TCR autorreativo). \\
\textbf{Memória virtual (\tc{VM})} &
Camund.: CD44\up{hi}, CD49d\up{low}, CD62L\up{hi}, CD122\up{hi} &
Células naive que proliferaram \emph{sem} reconhecimento de antígeno cognato, adquirindo
diferenciação parcial. \\
\end{longtable}
\normalsize

% =====================================================================
\section{Subpopulações de células T CD4\texorpdfstring{\up{+}}{+} (ênfase)}
\label{sec:cd4}
Recomendação editorial do consenso: usar o número do \emph{helper} em \textbf{subscrito}
(ex.: \tc{H}1, não ``TH1'' ou ``Th1''). Nas células CD4\up{+}, o ``H'' (de \emph{helper}) é,
na prática, redundante com CD4, mas foi mantido por estar consolidado.

\subsection{Células T auxiliares (\emph{helper}) e CD4\up{+} citotóxicas}
\renewcommand{\arraystretch}{1.3}
\small
\begin{longtable}{L{1.7cm} L{5.4cm} L{3.3cm} L{4.4cm}}
\caption{Subsets auxiliares CD4\up{+} e CD4\up{+} CTL.}\label{tab:helper}\\
\toprule
\textbf{Subset} & \textbf{Humano} & \textbf{Camundongo} & \textbf{Descrição / função} \\
\midrule
\endfirsthead
\toprule
\textbf{Subset} & \textbf{Humano} & \textbf{Camundongo} & \textbf{Descrição / função} \\
\midrule
\endhead
\bottomrule
\endfoot
\tc{H}1 &
CXCR3\up{+}, CCR4\up{--}, CCR6\up{--}, CXCR5\up{--}, T-bet\up{+}; \ifng{}\up{+} ao estímulo &
T-bet\up{+}; \ifng{}\up{+} ao estímulo &
Predispostas a \ifng{} (não IL-4/5/13). Dependem de \textbf{T-bet}. Imunidade a patógenos
intracelulares que escapam de macrófagos. \\
\tc{H}2 &
CRTh2\up{+}, CXCR3\up{--}, CCR4\up{+}, CCR6\up{--}, CXCR5\up{--}, GATA3\up{+};
IL-4/5/13\up{+} ao estímulo &
GATA3\up{+}; IL-4/5/13\up{+} ao estímulo &
Predispostas a IL-4, IL-5, IL-13 (não \ifng{}). Dependem de \textbf{GATA3}. Imunidade a
helmintos; inflamação alérgica. \\
\tc{H}9 &
CXCR3\up{--}, CCR4\up{+}, CCR6\up{--}, CXCR5\up{--}, CCR8\up{+}, PU.1\up{+}; IL-9 ao estímulo;
PPAR\ensuremath{\gamma}\up{+} na pele &
PU.1\up{+}; IL-9 ao estímulo &
Predispostas a IL-9. Dependem de \textbf{PU.1, IRF4, BATF}. Variantes: p\tc{H}2 /
\tc{H}2A (alérgica, humano) / \tc{H}9RM (memória residente, camundongo). \\
\tc{H}17 &
CXCR3\up{--}, CCR4\up{+/--}, CCR6\up{+}, CXCR5\up{--}, \rort{}\up{+}; IL-17\up{+} ao estímulo &
\rort{}\up{+}; IL-17\up{+} ao estímulo &
Predispostas a IL-17. Dependem de \textbf{\rort{}}. Imunidade de barreira (bactérias
extracelulares, fungos). Podem ser \emph{homeostáticas} ou \emph{patogênicas}. \\
CD4\up{+} CTL &
Ensaio de liberação de cromo; granzima~B\up{+}, CCR7\up{--}, CD27\up{--}; pode expressar
CD95 (FAS) & Igual ao humano &
Célula T CD4\up{+} com atividade citolítica. \\
\end{longtable}
\normalsize

\subsection{Família T folicular auxiliar (\tc{FH})}
As células \tc{FH} promovem ativação de células B, produção de anticorpos, hipermutação
somática, troca de classe e memória; dependem de \textbf{BCL6}. Podem ser subclassificadas
por citocinas (\tc{FH}1, \tc{FH}2). A definição de \tc{FR} e de \tc{FH} de centro germinativo
depende crucialmente da \emph{localização}, exigindo dados de imagem.

\renewcommand{\arraystretch}{1.3}
\small
\begin{longtable}{L{2.0cm} L{4.3cm} L{4.6cm} L{4.0cm}}
\caption{Subsets da família \tc{FH}.}\label{tab:tfh}\\
\toprule
\textbf{Subset} & \textbf{Humano} & \textbf{Camundongo} & \textbf{Descrição} \\
\midrule
\endfirsthead
\toprule
\textbf{Subset} & \textbf{Humano} & \textbf{Camundongo} & \textbf{Descrição} \\
\midrule
\endhead
\bottomrule
\endfoot
\tc{FH} & CXCR5\up{+}, BCL6\up{+} & Igual ao humano & Auxílio a células B; depende de BCL6. \\
GC \tc{FH} \newline (centro germinativo) & ND &
BCL6\up{hi}, CXCR5\up{hi}, PD-1\up{hi}, ICOS\up{+}, FR4\up{hi}, PSGL1\up{low} &
Residem em centros germinativos; selecionam células B de alta afinidade. \\
MZ \tc{FH} \newline (zona do manto) & ND &
BCL6\up{low}, CXCR5\up{low}, PD-1\up{low}, FR4\up{low}, PSGL1\up{hi} &
Na zona do manto; possíveis precursoras de GC \tc{FH} (ou GC \tc{FH} que retornaram). \\
\tc{FHM} \newline (memória) & CXCR4\up{+}, FOXP3\up{--} &
TCF1\up{+}, FR4\up{+}, P2RX7\up{+}, CXCR4\up{+} &
\tc{FH} longevas que persistem, supostamente, após a eliminação do antígeno. \\
\end{longtable}
\normalsize

\subsection{Família reguladora (\tc{reg})}
As células \tc{reg} regulam respostas imunes e promovem tolerância; dependem de
\textbf{FOXP3}. Marcadores comuns (humano e camundongo): FOXP3\up{+}, CD25\up{hi},
CD127\up{low}, CTLA4\up{hi}, GITR\up{hi}. \emph{Nota importante:} em humanos não há
marcadores estabelecidos para distinguir t\tc{reg} de p\tc{reg}; no camundongo,
neuropilina-1 e Helios fazem essa distinção.

\renewcommand{\arraystretch}{1.3}
\small
\begin{longtable}{L{1.9cm} L{4.6cm} L{4.8cm} L{3.6cm}}
\caption{Subsets de células T reguladoras.}\label{tab:treg}\\
\toprule
\textbf{Subset} & \textbf{Humano} & \textbf{Camundongo} & \textbf{Descrição} \\
\midrule
\endfirsthead
\toprule
\textbf{Subset} & \textbf{Humano} & \textbf{Camundongo} & \textbf{Descrição} \\
\midrule
\endhead
\bottomrule
\endfoot
\tc{reg} & FOXP3\up{+}, CD25\up{hi}, CD127\up{low}, CTLA4\up{+}, GITR\up{+} & Igual ao humano &
Regulação imune e tolerância; depende de FOXP3. \\
t\tc{reg} \newline (tímica / natural) &
FOXP3\up{+}, CCR7\up{hi}, CD45RA\up{+}, CD45RO\up{--} &
FOXP3\up{+}, CD62L\up{+}, CD44\up{low}, neuropilina-1\up{hi}, Helios\up{hi} &
Diferenciam-se por autoantígenos no timo. Difíceis de distinguir de p\tc{reg} em humanos. \\
p\tc{reg} \newline (periférica / induzida) &
FOXP3\up{+}, CCR7\up{hi}, CD45RA\up{+}, CD45RO\up{--} &
FOXP3\up{+}, CD62L\up{+}, CD44\up{low}, neuropilina-1\up{low}, Helios\up{low/int}
(\rort{}\up{+} ou GATA3\up{hi}) &
Diferenciam-se por antígenos (self, comensais, dieta) na periferia. \\
e\tc{reg} \newline (efetora) &
FOXP3\up{hi}, CD45RO\up{+}, CCR8\up{hi}, ICOS\up{hi}, IRF4\up{+} &
FOXP3\up{hi}, CD62L\up{--}, CD44\up{hi}, BLIMP1\up{+}, IL-10\up{+}, TGF\ensuremath{\beta}\up{+} &
t\tc{reg}/p\tc{reg} com diferenciação adicional; maior secreção de citocinas supressoras
(IL-10, TGF\ensuremath{\beta}). \\
i\tc{reg} \newline (\emph{in vitro}) & ND & ND &
CD4\up{+} induzidas a \tc{reg} \emph{in vitro} por citocinas. \\
\tc{FR} \newline (folicular reg.) &
CXCR5\up{+}, FOXP3\up{+}, GITR\up{+}, IL-10\up{+} &
CXCR5\up{+}, PD-1\up{+}, BCL6\up{+}, FOXP3\up{+}, GITR\up{+}, IL-10\up{+} &
\tc{reg} localizadas em folículos B / centros germinativos (classificação requer imagem). \\
\end{longtable}
\normalsize

% =====================================================================
\section{Subsets de memória (CD4\texorpdfstring{\up{+}}{+} ou CD8\texorpdfstring{\up{+}}{+})}
Os termos abaixo são usados independentemente da especialização funcional (ex.: \tc{H}1 vs.\ \tc{H}2).

\renewcommand{\arraystretch}{1.25}
\small
\begin{longtable}{L{2.2cm} L{6.2cm} L{6.6cm}}
\caption{Subsets de memória.}\label{tab:mem}\\
\toprule
\textbf{Subset} & \textbf{Marcadores (proxies)} & \textbf{Descrição} \\
\midrule
\endfirsthead
\toprule
\textbf{Subset} & \textbf{Marcadores (proxies)} & \textbf{Descrição} \\
\midrule
\endhead
\bottomrule
\endfoot
\tc{CM} (memória central) &
Humano: CD62L\up{+}, CCR7\up{+}. Camund.: CD44\up{hi}, CD62L\up{+}, TCF1\up{+}, CCR7\up{+} &
Recirculam entre linfa e sangue; cruzam HEVs em repouso; mantêm capacidade de
proliferar/diferenciar. \\
\tc{EM} (memória efetora) &
CCR7\up{--}, CD45RA\up{--} (humano); CD44\up{hi}, CD62L\up{--} (camund.) &
Carecem de receptores de \emph{homing} (CD62L, CCR7). Termo ``guarda-chuva'' para ampla gama
de estados. \\
\tc{EMRA} &
CCR7\up{--}, CD45RA\up{+}; CD27\up{--}, CX3CR1\up{hi}, CD57\up{+} &
Subset humano de \tc{EM} que re-expressa CD45RA. \textbf{Não se aplica a camundongos.} \\
\tc{RM} (residente) &
KLF2\up{low}, CD69\up{+}, CD62L\up{--}, CCR7\up{--}; freq.\ CD103\up{+}, CD49a\up{+}, CXCR6\up{+} &
Residem em um único tecido/órgão sem recircular pelo sangue. Marcadores variam por órgão. \\
Ex\tc{RM} &
Fenocopia \tc{EM}/\tc{CM}; pode reter CD43 elevado (clone 1B11) &
\tc{RM} que se diferenciou e voltou à circulação; pode reter imprint epigenético de residência. \\
\tc{PM} (memória periférica) &
CD45RA\up{--}, CCR7\up{int}, CD27\up{+}; CX3CR1\up{int} (camund.) &
Recirculam entre tecidos não linfoides e sangue (sem marcadores definitivos). \\
\tc{SCM} (memória tronco) &
CD95\up{+}, CCR7\up{+}, CD27\up{+}, CD28\up{+}, CD45RA\up{+}, Eomes\up{--} &
Raras, pouco diferenciadas; auto-renovação + diferenciação em outros subsets.
Compartilham marcadores com naive. \\
\end{longtable}
\normalsize

% =====================================================================
\section{Subsets de exaustão e efetores CD8\texorpdfstring{\up{+}}{+} (síntese)}
Embora centrados em CD8\up{+}, são incluídos para completude.

\noindent\textbf{Efetores/ativados CD8\up{+} (Tabela~2):} \tc{SLEC}/TE (efetor terminal,
citotóxica, morre na transição --- KLRG1\up{+}, CD127\up{low}); \tc{MPEC} (precursora de memória,
sobrevive à transição --- CD127\up{+}, CD27\up{+}, TCF1\up{+}); \tc{EEC} (efetora inicial ---
KLRG1\up{low}, CXCR3\up{+}, produz \ifng{}/TNF); \tc{LLEC} (efetora longeva, muito citotóxica,
prolifera mal --- KLRG1\up{hi}, CX3CR1\up{hi}, granzima~B\up{hi}).

\vspace{4pt}
\noindent\textbf{Exaustão (Tabela~5)} --- os marcadores variam por contexto (câncer vs.\
infecção), sobretudo em humanos:
\begin{itemize}[leftmargin=1.4em,itemsep=2pt]
  \item \textbf{\tc{PEX} (progenitora/precursora exausta):} sustenta a resposta na estimulação
        crônica, auto-renova e responde ao bloqueio de PD-1. PD-1\up{+}, TOX\up{+}, TCF1\up{+},
        SLAMF6\up{+}, CXCR5\up{+}, TIM3\up{--}. Também chamada ``stem-like''. Módulo modular: \textbf{X\textsubscript{p}}.
  \item \textbf{\tc{EX-int} / \tc{EX-eff} (intermediária / efetora-símile):} transicional;
        expressa granzima~B, perforina, \ifng{}; PD-1\up{+}, TIM3\up{+}, T-bet\up{+}, TCF1\up{--}.
  \item \textbf{\tc{EX-term} (terminalmente exausta):} baixa proliferação, função efetora
        reduzida; PD-1\up{+}, TOX\up{+}, TIM3\up{+}, CD39\up{+}, CD101\up{+}, TCF1\up{--}. Em câncer,
        frequentemente chamada ``disfuncional''. Módulo modular: \textbf{X\textsubscript{t}}.
\end{itemize}

\noindent\textbf{Innate-like (Tabela~6, menção):} MAIT (TCR$\alpha$ invariante, MR1),
iNKT (CD1d-$\alpha$GalCer), \gd{} T (definidas pelo TCR\gd{}; fortemente associadas ao uso
V$\gamma$/V$\delta$), CD8$\alpha\alpha$\up{+} TCR$\alpha\beta$, duplo-negativas (DN) e
duplo-positivas extratímicas (DP).

% =====================================================================
\section{Nomenclatura modular v1.0}
\label{sec:modular}
\textbf{Ideia central:} assim como 26 letras geram $\sim$1 milhão de palavras, poucos
descritores combináveis capturam a complexidade das células T \emph{sem} inventar rótulos
novos. O esquema descreve propriedades individuais e \textbf{não obriga} o pesquisador a
inferir o que não mediu (permite ``desconhecido''). Evita-se ``C'' e ``E'' para não confundir
com \tc{CM}/\tc{EM}.

\renewcommand{\arraystretch}{1.25}
\small
\begin{longtable}{L{3.2cm} L{2.0cm} L{9.6cm}}
\caption{Módulos da nomenclatura modular v1.0 (Fig.~1 do artigo).}\label{tab:mod}\\
\toprule
\textbf{Módulo} & \textbf{Código} & \textbf{Significado} \\
\midrule
\endfirsthead
\toprule
\textbf{Módulo} & \textbf{Código} & \textbf{Significado} \\
\midrule
\endhead
\bottomrule
\endfoot
Linhagem & CD4\up{+}, CD8\up{+}, \gd{}\up{+}, NKT, MAIT, CD8$\alpha\alpha$\up{+} &
Pode estender-se a NK, ILC. \\
Função & \tc{H}1, \tc{reg}, \tc{C}2\ldots & Subsets funcionais atuais (idem prática vigente). \\
\midrule
\multirow{3}{*}{Migração}
 & \textbf{S} & \emph{Secondary lymphoid organ}: entra em SLO não inflamado a partir do sangue
                (proxy: CD62L\up{+} e/ou CCR7\up{+}). \\
 & \textbf{D} & \emph{Disseminated}: não entra em SLO não inflamado (proxy: CD62L\up{--} e/ou
                CCR7\up{--}). Engloba \tc{EM}, \tc{PM} e \tc{RM}. \\
 & \textbf{U} & \emph{Unknown}: propriedades de migração desconhecidas / não avaliadas. \\
\midrule
\multirow{3}{*}{\shortstack[l]{Migração\\(subscrito)}}
 & \textsubscript{B} & \emph{Blood}: isolada do sangue, sem outras alegações (ex.: D\textsubscript{B}). \\
 & \textsubscript{W} & \emph{Widespread}: recircula por órgãos não linfoides (ex.: D\textsubscript{W}). \\
 & \textsubscript{R} & \emph{Resident}: estacionada em órgãos/compartimentos (ex.: D\textsubscript{R}). \\
\midrule
\multirow{5}{*}{\shortstack[l]{Estado de\\diferenciação}}
 & \textbf{N} & Naive. \\
 & \textbf{A} & Ativada. \\
 & \textbf{M} & Memória. \\
 & \textbf{X} & Exausta. \\
 & \textbf{G} & Anérgica. \\
\midrule
\shortstack[l]{Estado\\(subscrito)}
 & \textsubscript{p}, \textsubscript{t} & \emph{progenitor} / \emph{terminal}
   (ex.: \tc{PEX}=X\textsubscript{p}; \tc{SCM}=M\textsubscript{p}; MPEC=A\textsubscript{p}; SLEC=A\textsubscript{t}). \\
\midrule
\multirow{2}{*}{Antígeno}
 & \textbf{+} & Antígeno presente no organismo (ex.: M+). \\
 & \textbf{0} & Antígeno eliminado ou irrelevante para a célula (ex.: M0). \\
\end{longtable}
\normalsize

\noindent\emph{Ordem de leitura:} \texttt{[Linhagem] [Função] T[Migração][subscrito
migração][Estado][subscrito estado][Antígeno]}. Ex.: \textbf{CD4\up{+} \tc{H}2D\textsubscript{R}M0}.

\subsection{Exemplos aplicados a CD4\up{+} (Fig.~3 e Tabela~7)}
\renewcommand{\arraystretch}{1.3}
\small
\begin{longtable}{L{8.6cm} L{6.0cm}}
\caption{Descrição em prosa vs.\ nomenclatura modular (exemplos CD4\up{+}).}\\
\toprule
\textbf{Descrição da célula} & \textbf{Modular} \\
\midrule
\endfirsthead
\toprule
\textbf{Descrição da célula} & \textbf{Modular} \\
\midrule
\endhead
\bottomrule
\endfoot
CD4\up{+} sem outra caracterização & CD4\up{+} T cell \\
\tc{H}2 CD62L\up{--}, recém-estimulada, isolada do sangue & CD4\up{+} \tc{H}2D\textsubscript{B}A \\
CD4\up{+} do pulmão, CCR7\up{--}, CD69\up{+}, migração desconhecida & Pulmão CD4\up{+} TD \\
\ldots{} idem, alegada recém-ativada & Pulmão CD4\up{+} TDA \\
\ldots{} idem, alegada memória residente & Pulmão CD4\up{+} TD\textsubscript{R}M \\
CD4\up{+} \tc{H}17 do pulmão, residente, antígeno eliminado &
Pulmão CD4\up{+} \tc{H}17D\textsubscript{R}M0 \\
\midrule
\multicolumn{2}{l}{\textit{Progressão de detalhe para a mesma célula \tc{H}2 (Fig.~3):}}\\
Só linhagem $\rightarrow$ + função $\rightarrow$ + memória $\rightarrow$ + migração $\rightarrow$ + antígeno &
CD4\up{+} $\rightarrow$ CD4\up{+}\tc{H}2 $\rightarrow$ CD4\up{+}\tc{H}2\,M $\rightarrow$
CD4\up{+}\tc{H}2D\,M $\rightarrow$ CD4\up{+}\tc{H}2D\textsubscript{R}M0 \\
\end{longtable}
\normalsize

\noindent\textbf{Observação sobre função de CD8\up{+}:} pode-se usar ``C'' (historicamente
``citotóxica'', hoje só ``CD8\up{+}''): p.ex.\ CD8\up{+}\,\tc{C}2, \tc{C}17, \tc{C}22, ou
CD8\up{+}\,\tc{reg}. Células citotóxicas CD4\up{+} \emph{ou} CD8\up{+} podem ser \tc{CTL}.

% =====================================================================
\section{Conclusões}
\begin{itemize}[leftmargin=1.4em,itemsep=2pt]
  \item As diretrizes entregam: (i) definições padronizadas de termos existentes;
        (ii) o pedido de explicitar critérios nos Métodos; (iii) a nomenclatura modular.
  \item A nomenclatura modular \textbf{separa} função, migração, estado de diferenciação e
        status de antígeno, dando precisão sob demanda \emph{sem} forçar alegações não medidas
        e \emph{sem} inflar o vocabulário.
  \item É extensível (fatores de transcrição, epigenética) e pode servir também a NK e ILC.
        Deve coexistir com a nomenclatura clássica por algum tempo; os autores podem usar uma,
        outra, ou ambas no mesmo artigo, desde que a comunicação seja inequívoca.
\end{itemize}

\vfill
\begin{center}\small\color{softgray}
Pôster oficial das diretrizes:
\href{https://doi.org/10.1038/s41577-025-01246-2}{doi:10.1038/s41577-025-01246-2}.
\end{center}

\end{document}
